% GENERATED FILE. Edit canonical lessons, then run npm run generate:books.
% canonical-source-hash: fnv1a64:ddf9837f61b8f332
% canonical-lessons: ZH-C480-touyun, ZH-C480-fayang, ZH-C480-liubiti, ZH-C480-guomin, ZH-C480-shimian, ZH-R480-review-476-480

\chapter{Symptoms}
\label{ch:zh-symptoms}
\hlchaptermodality{480}

\begin{chapteropening}
\textbf{By the end of this chapter:} \emph{I can say I feel dizzy or itchy, that my nose is running, that I am allergic to something, and that I cannot sleep.}

Complete the last lesson of chapter 480: I can say I feel dizzy or itchy, that my nose is running, that I am allergic to something, and that I cannot sleep.
\end{chapteropening}

\section[tóuyūn]{tóuyūn — dizzy}

\label{lesson:ZH-C480-touyun}



\begin{quote}
Before the new one: say the Mandarin for to give blood, then the Mandarin for blood pressure.
\end{quote}

\subsection*{You'll want to know: tóuyūn}
\emph{tóuyūn} — \textquotedblleft{}dizzy\textquotedblright{}.

\begin{quote}
\emph{wǒ gǎnjué yǒudiǎnr tóuyūn} — I feel a bit dizzy.
\end{quote}

Its characters come later; for now it is a word for the ear.

\begin{cousinweb}
\emph{tóuyūn} is built from \zh{头} 'head' and \emph{yūn}, both inherited from Middle Chinese, \emph{duw} and \emph{hjunH} (Guangyun, in Baxter's transcription), which glosses the second as a halo round the sun or moon. The level tone of \emph{yūn} is irregular, and the sense 'dizzy' is not given there.
\end{cousinweb}

\begin{grammarlens}[title={What you've built}]
One more word for saying how it was.
\end{grammarlens}

\subsection*{Your turn}
\begin{itemize}
\raggedright
  \item \emph{Say it:} \emph{tóuyūn}
  \item \emph{Say it:} \emph{tóuyūn}, once more
  \item \emph{Say it:} \emph{wǒ gǎnjué yǒudiǎnr tóuyūn} — I feel a bit dizzy.
  \item \emph{Recall:} say the Mandarin for to give blood, then the Mandarin for blood pressure, then say \emph{tóuyūn} again
\end{itemize}

\begin{tcolorbox}[breakable,colback=teal!4,colframe=teal!35!black,title={Before you move on}]
What does tóuyūn mean? (\textquotedblleft{}Dizzy\textquotedblright{}.) Say it once more.
\end{tcolorbox}
\section[fāyǎng]{fāyǎng — to itch}

\label{lesson:ZH-C480-fayang}



\begin{quote}
Before the new one: say the Mandarin for blood pressure, then the Mandarin for dizzy.
\end{quote}

\subsection*{You'll want to know: fāyǎng}
\emph{fāyǎng} — \textquotedblleft{}to itch\textquotedblright{}.

\begin{quote}
\emph{shāngkǒu yǒudiǎnr fāyǎng} — The wound is a bit itchy.
\end{quote}

Its characters come later; for now it is a word for the ear.

\begin{cousinweb}
\emph{fāyǎng} is built from \zh{发} 'to set going' and \emph{yǎng} 'itchy', both inherited from Middle Chinese, \emph{pjot} and \emph{yangX} (Guangyun, in Baxter's transcription), which glosses the second as itchy skin.
\end{cousinweb}

\begin{grammarlens}[title={What you've built}]
One more word for saying how it was.
\end{grammarlens}

\subsection*{Your turn}
\begin{itemize}
\raggedright
  \item \emph{Say it:} \emph{fāyǎng}
  \item \emph{Say it:} \emph{fāyǎng}, once more
  \item \emph{Say it:} \emph{shāngkǒu yǒudiǎnr fāyǎng} — The wound is a bit itchy.
  \item \emph{Recall:} say the Mandarin for blood pressure, then the Mandarin for dizzy, then say \emph{fāyǎng} again
\end{itemize}

\begin{tcolorbox}[breakable,colback=teal!4,colframe=teal!35!black,title={Before you move on}]
What does fāyǎng mean? (\textquotedblleft{}To itch\textquotedblright{}.) Say it once more.
\end{tcolorbox}
\section[liú bítì]{liú bítì — to have a runny nose}

\label{lesson:ZH-C480-liubiti}



\begin{quote}
Before the new one: say the Mandarin for dizzy, then the Mandarin for to itch.
\end{quote}

\subsection*{You'll want to know: liú bítì}
\emph{liú bítì} — \textquotedblleft{}to have a runny nose\textquotedblright{}.

\begin{quote}
\emph{wǒ gǎnmào le liǎng tiān xiànzài hái zài liú bítì} — I've had a cold for two days, and my nose is still running.
\end{quote}

Its characters come later; for now it is a word for the ear.

\begin{cousinweb}
\emph{liú bítì} is built from \emph{liú} 'to flow', \emph{bí} 'nose' and \emph{tì}, all inherited from Middle Chinese, \emph{ljuw}, \emph{bjijH} and \emph{thejH} (Guangyun, in Baxter's transcription). The dictionary explains the first, after an older book, as water moving, and glosses the third as tears; the sense 'mucus of the nose' is not given there.
\end{cousinweb}

\begin{grammarlens}[title={What you've built}]
One more word for saying how it was.
\end{grammarlens}

\subsection*{Your turn}
\begin{itemize}
\raggedright
  \item \emph{Say it:} \emph{liú bítì}
  \item \emph{Say it:} \emph{liú bítì}, once more
  \item \emph{Say it:} \emph{wǒ gǎnmào le liǎng tiān xiànzài hái zài liú bítì} — I've had a cold for two days, and my nose is still running.
  \item \emph{Recall:} say the Mandarin for dizzy, then the Mandarin for to itch, then say \emph{liú bítì} again
\end{itemize}

\begin{tcolorbox}[breakable,colback=teal!4,colframe=teal!35!black,title={Before you move on}]
What does liú bítì mean? (\textquotedblleft{}To have a runny nose\textquotedblright{}.) Say it once more.
\end{tcolorbox}
\section[guòmǐn]{guòmǐn — to be allergic}

\label{lesson:ZH-C480-guomin}



\begin{quote}
Before the new one: say the Mandarin for to itch, then the Mandarin for to have a runny nose.
\end{quote}

\subsection*{You'll want to know: guòmǐn}
\emph{guòmǐn} — \textquotedblleft{}to be allergic\textquotedblright{}.

\begin{quote}
\emph{wǒ duì huā guòmǐn} — I'm allergic to flowers.
\end{quote}

Its characters come later; for now it is a word for the ear.

\begin{cousinweb}
\emph{guòmǐn} is built from \zh{过} 'beyond' and \emph{mǐn} 'quick', both inherited from Middle Chinese, \emph{kwaH} and \emph{minX} (Guangyun, in Baxter's transcription), which glosses the first 'to go beyond' and the second 'quick' and 'clever': too quick to react. The medical sense is modern.
\end{cousinweb}

\begin{grammarlens}[title={What you've built}]
One more word for saying how it was.
\end{grammarlens}

\subsection*{Your turn}
\begin{itemize}
\raggedright
  \item \emph{Say it:} \emph{guòmǐn}
  \item \emph{Say it:} \emph{guòmǐn}, once more
  \item \emph{Say it:} \emph{wǒ duì huā guòmǐn} — I'm allergic to flowers.
  \item \emph{Recall:} say the Mandarin for to itch, then the Mandarin for to have a runny nose, then say \emph{guòmǐn} again
\end{itemize}

\begin{tcolorbox}[breakable,colback=teal!4,colframe=teal!35!black,title={Before you move on}]
What does guòmǐn mean? (\textquotedblleft{}To be allergic\textquotedblright{}.) Say it once more.
\end{tcolorbox}
\section[shīmián]{shīmián — to be unable to sleep; to lie awake at night}

\label{lesson:ZH-C480-shimian}



\begin{quote}
Before the new one: say the Mandarin for to have a runny nose, then the Mandarin for to be allergic.
\end{quote}

\subsection*{You'll want to know: shīmián}
\emph{shīmián} — \textquotedblleft{}to be unable to sleep; to lie awake at night\textquotedblright{}.

\begin{quote}
\emph{wǒ shīmián le yí ge xīngqī} — I haven't been able to sleep for a week.
\end{quote}

Its characters come later; for now it is a word for the ear.

\begin{cousinweb}
\emph{shīmián} is built from \emph{shī} 'to lose' and \emph{mián} 'sleep', both inherited from Middle Chinese, \emph{syit} and \emph{men} (Guangyun, in Baxter's transcription), which glosses the first 'to let slip' and the second 'to sleep'.
\end{cousinweb}

\begin{grammarlens}[title={What you've built}]
One more word for saying how it was.
\end{grammarlens}

\subsection*{Your turn}
\begin{itemize}
\raggedright
  \item \emph{Say it:} \emph{shīmián}
  \item \emph{Say it:} \emph{shīmián}, once more
  \item \emph{Say it:} \emph{wǒ shīmián le yí ge xīngqī} — I haven't been able to sleep for a week.
  \item \emph{Recall:} say the Mandarin for to have a runny nose, then the Mandarin for to be allergic, then say \emph{shīmián} again
\end{itemize}

\begin{tcolorbox}[breakable,colback=teal!4,colframe=teal!35!black,title={Before you move on}]
What does shīmián mean? (\textquotedblleft{}To be unable to sleep; to lie awake at night\textquotedblright{}.) Say it once more.
\end{tcolorbox}
\section[(dialogue)]{Looking back: from Neck and Shoulders to Symptoms}

\label{lesson:ZH-R480-review-476-480}



\begin{quote}
Say the Mandarin for the neck and the tongue.
\end{quote}

\subsection*{Your turn}
\begin{itemize}
\raggedright
  \item \emph{Say it:} the neck, the tongue, the shoulder, to shake one's head, bodily pain
  \item \emph{Say it:} the heart as an organ, a bone, the brain, one's digestion, an appetite
  \item \emph{Say it:} to scald oneself, an open wound, badly hurt, to sprain, to dress a wound
  \item \emph{Say it:} a scar, blood, to stop the bleeding, to give blood, blood pressure
  \item \emph{Say it:} dizzy, to itch, to have a runny nose, to be allergic, to lie awake at night
  \item \emph{Read it:} \zh{伤} \zh{血} — each one aloud, then a word that holds it
\end{itemize}

\begin{tcolorbox}[breakable,colback=teal!4,colframe=teal!35!black,title={Before you move on}]
The neck? (\emph{bózi}.) The tongue? (\emph{shétou}.) The shoulder? (\emph{jiānbǎng}.) To shake one's head? (\emph{yáotóu}.) Bodily pain? (\emph{téngtòng}.) The heart as an organ? (\emph{xīnzàng}.) A bone? (\emph{gǔtou}.) The brain? (\emph{dànǎo}.) One's digestion? (\emph{chángwèi}.) An appetite? (\emph{wèikǒu}.) To scald oneself? (\emph{tàngshāng}.) An open wound? (\textbf{\zh{伤口}}, \emph{shāngkǒu}.) Badly hurt? (\textbf{\zh{重伤}}, \emph{zhòngshāng}.) To sprain? (\emph{niǔshāng}.) To dress a wound? (\emph{bāozā}.) A scar? (\emph{shāngbā}.) Blood? (\textbf{\zh{血}}, \emph{xiě}.) To stop the bleeding? (\emph{zhǐxuè}.) To give blood? (\emph{xiànxuè}.) Blood pressure? (\emph{xuèyā}.) Dizzy? (\emph{tóuyūn}.) To itch? (\emph{fāyǎng}.) To have a runny nose? (\emph{liú bítì}.) To be allergic? (\emph{guòmǐn}.) To lie awake at night? (\emph{shīmián}.)
\end{tcolorbox}
